%%%%%%%%%%%%%%%%%%%%%%%%%%%%%%%%%%%%%%%%%
% Medium Length Professional CV
% LaTeX Template
% Version 2.0 (8/5/13)
%
% This template has been downloaded from:
% http://www.LaTeXTemplates.com
%
% Original author:
% Trey Hunner (http://www.treyhunner.com/)
%
% Important note:
% This template requires the resume.cls file to be in the same directory as the
% .tex file. The resume.cls file provides the resume style used for structuring the
% document.
%
%%%%%%%%%%%%%%%%%%%%%%%%%%%%%%%%%%%%%%%%%

%----------------------------------------------------------------------------------------
%	PACKAGES AND OTHER DOCUMENT CONFIGURATIONS
%----------------------------------------------------------------------------------------

\documentclass{resume} % Use the custom resume.cls style

\usepackage[left=0.4 in,top=0.4in,right=0.4 in,bottom=0.4in]{geometry} % Document margins
\usepackage{fontawesome}
\usepackage{setspace}
\usepackage{hyperref}
\newcommand{\tab}[1]{\hspace{.2667\textwidth}\rlap{#1}} 
\newcommand{\itab}[1]{\hspace{0em}\rlap{#1}}
\name{Gahan Saraiya} % Your name
%\address{123 Pleasant Lane \\ City, State 12345} % Your secondary addess (optional)
\address{ 
  \faPhone \ +91 7359508173 \\
  \faEnvelope \ \href{mailto: gahansaraiya@gmail.com}{gahansaraiya@gmail.com} \\  
  \faLinkedin \ \href{https://www.linkedin.com/in/gahan-saraiya/}{linkedin.com/in/gahan-saraiya} \\ \faGithub \   \href{https://github.com/gahan9/}{github.com/gahan9} 
}  % Your phone number and email

\hypersetup{
    colorlinks=true,
    linkcolor=blue,
    filecolor=magenta,      
    urlcolor=blue,
    pdfpagemode=FullScreen,
    }
\begin{document}

% AI/ML - 0
% DevOps - 1
% PhD - 2
\def\profile{0}

%----------------------------------------------------------------------------------------
%	OBJECTIVE
%----------------------------------------------------------------------------------------

\begin{rSection}{OBJECTIVE}
{High Performance Computing (HPC) Development Engineer with 8+ years of Research and Development experience in system software, GPU computing, and distributed systems architecture. Currently at AMD India Pvt Ltd., specializing in performance optimization and AI/ML workload execution on HPC accelerators (GPUs). Deep expertise in Python, modern C++, and Bash scripting for large-scale on-prem and cloud infrastructure deployments. Strong background in operating systems internals, computer networks, and high-performance applications across heterogeneous computing environments comprising CPUs, GPUs, and accelerators. Proven track record in deploying and optimizing Deep-Learning frameworks including TensorFlow and PyTorch on enterprise HPC clusters. Passionate about open-source contribution and advancing the state-of-the-art in parallel computing and distributed systems.}

\end{rSection}




%----------------------------------------------------------------------------------------
%	TECHNICAL STRENGTHS SECTION
%----------------------------------------------------------------------------------------

\begin{rSection}{TECHNICAL SKILLS}

\begin{tabular}{ @{} >{\bfseries}l @{\hspace{6ex}} l }

% PROFILE 0
\if\profile0
Programming Languages & \textbf{C++}, C, \textbf{Python}, \textbf{Bash/Shell Scripting}, SQL
\\
HPC \& GPU Computing & ROCm, \textbf{HIP Programming}, \textbf{CUDA}, PyCUDA, OpenCL, MPI, OpenMP,
\\ & GPU Profiling, Memory Optimization, Kernel Development
\\
Deep Learning Frameworks & \textbf{TensorFlow}, \textbf{PyTorch}, TensorRT, ONNX Runtime,
\\ & Large-Scale Model Training, Distributed Training
\\
Distributed Systems & CPU/GPU/Accelerator Architecture, Parallel Programming,
\\ & High-Performance Networks, InfiniBand, RDMA, MPI
\\
Operating Systems & Linux Internals, Kernel Development, UEFI/BIOS Firmware,
\\ & System Performance Tuning, Memory Management
\\
Cloud \& Infrastructure & AWS (EC2, S3, SageMaker), GCP, Azure, Kubernetes, Docker,
\\ & Terraform, Large-Scale On-Prem HPC Clusters, Cloud Security Policies
\\
CI/CD \& DevOps & Jenkins, GitHub Actions, TeamCity, JFrog Artifactory,
\\ & Cluster Validation CI
\\
System Level IP & RAS (Reliability, Availability, Serviceability), Compute Partitioning,
\\ & Boot-time Optimization, CPER, SMI
\\
\fi

%  PROFILE 1
\if\profile1
Languages & Python, Shell Scripting, SQL, LaTeX, YAML 
\\
Frameworks & Pytorch, Pytest, Edk2, PyCUDA, Pandas, UEFI, Selenium, Redfish
\\
Architecture Design & Parallel Programming System Architecture, ROCm, HIP Programming
\\
Tools and Services & GitHub Actions, Jenkins, Restful APIs, Teamcity, Terraform,
\\ & Docker, Docker Compose, Jenkins, TeamCity, GiT 
\\
Cloud Technologies & AWS EC2, AWS S3, AWS Cloudwatch, VPC Networking, AWS Sage Maker, 
\\ & AWS Lambda, AWS Cloudfront, ECR, ECS, EKS, JFrog Artifactory,
\\ & GCP, Azure, Digital Ocean, Cloud Security Policies, Cluster Validation CI
    
\fi

\end{tabular}

\end{rSection}


%----------------------------------------------------------------------------------------

%----------------------------------------------------------------------------------------
%	Experience SECTION
%----------------------------------------------------------------------------------------

\begin{rSection}{Experience}

{\bf Senior System Software Designer at AMD India Pvt Ltd} \hfill {Mar 2024 - Present }
\begin{itemize}
\if\profile1
    \item Workload execution pool management for cross geo platforms for CI/CD on Jenkins
    \item Release management of ROCm builds for production, non-production and various other categories with matrix of configuration of OS and toolchain
    \item Docker orchestration on enterprise registry
    \item Platform provisioning migration with terraform
\fi
    \item \textbf{R\&D Leadership}: Spearheading research and development of high-performance computing solutions for next-generation AMD GPU accelerators (MI2XX, MI3XX series).
    \item \textbf{Performance Engineering}: Architecting and implementing comprehensive GPU benchmarking frameworks using \textbf{modern C++} and \textbf{Python} for functional and performance validation of HPC devices.
    \item \textbf{Deep Learning Framework Optimization}: Deploying and optimizing \textbf{PyTorch} workloads on large-scale GPU clusters, implementing GEMM operations and distributed training pipelines.
    \item \textbf{Distributed Systems Architecture}: Designing workload orchestration systems across heterogeneous computing environments comprising \textbf{CPUs, GPUs, and accelerators}.
    \item \textbf{System-Level Optimization}: Developing load-balanced GPU performance criteria with thermal management, power consumption profiling, memory optimization, and core utilization metrics.
    \item \textbf{Operating Systems Expertise}: Implementing RAS (Reliability, Availability, Serviceability) injection mechanisms and System Management Interface (SMI) profiling.
    \item \textbf{HIP/ROCm Development}: Designing firmware capabilities including compute partitioning, HIP kernel profiling, and parallel execution strategies.
\if\profile0
    \item Collaborating with cross-functional R\&D teams of engineers and researchers on cutting-edge HPC solutions.
\fi
\end{itemize}

{\textbf{Senior Firmware Development Engineer at Intel Technology India Pvt Ltd}}  \hfill {Oct 2023 - Mar 2024}
\begin{itemize}
    \item \textbf{HPC Infrastructure R\&D}: Platform provisioning and infrastructure design for mass automation of high-performance computing products.
    \item \textbf{System Architecture Design}: Led Platform Orchestration Layer architecture for Universal Scalable Firmware (USF) using \textbf{modern C++} and \textbf{Python}.
    \item \textbf{AI/ML Integration}: Implemented Azure OpenAI Generative AI solutions for automated code review and security analysis.
\if\profile1
    \item Terraform deployments for environment for on-the-go dev environment over kubernetes cluster. (including setting up a node platform from scratch)
\fi
    \item \textbf{Large-Scale Infrastructure}: Designed CI/CD pipelines with GitHub Actions and Jenkins for distributed build systems across baremetal and containerized environments.
    \item \textbf{Operating System Optimization}: Enhanced firmware-based boot-time performance through deep OS-level optimizations.
\end{itemize}
 
{\textbf{Firmware Development Engineer at Intel Technology India Pvt Ltd}}  \hfill {Jun 2020 - Sept 2023}
\begin{itemize}
    \item \textbf{R\&D Technical Lead}: Spearheaded next-generation platform architecture design, owning domains including Automation Framework, Platform Provisioning, and Quantum Firmware Design.
    \item \textbf{C++ and Python Development}: Developed high-performance firmware validation tools using \textbf{Python scripting} for large-scale automated testing.
    \item \textbf{Distributed Systems}: Managed end-to-end SDLC for distributed firmware systems including security analysis (Snyk, Coverity) and static code analysis.
    \item \textbf{Scripting \& Automation}: Developed \textbf{Bash} and \textbf{Python} scripts for automated testing pipelines, GitHub hooks, and CI/CD integration with Jenkins/TeamCity.
    \item \textbf{Operating Systems Expertise}: Deep involvement in UEFI/BIOS firmware development and system-level debugging.
    \item Mentoring junior engineers on HPC concepts, parallel programming, and modern C++ best practices.
\end{itemize}

{\textbf{Firmware Engineer Intern at Intel Technology India Pvt Ltd}}  \hfill {Jun 2019 - May 2020}
\begin{itemize}
    \item \textbf{Low-Level Systems Development}: BIOS/UEFI Firmware development using EDK2 framework with \textbf{C} and \textbf{C++} for client products.
    \item \textbf{Operating Systems}: Deep understanding of boot processes, memory management, and hardware abstraction layers.
\end{itemize}

% \if\profile2
{\textbf{Teaching Assistantship at Nirma University}}  \hfill {Aug 2018 - May 2019}
\begin{itemize}
    \item Data Mining and Visulaization
\end{itemize}
% \fi

% \if\profile1
{\textbf{Software Developer at Quixom Technology}}   \hfill {Jul 2017 - Apr 2018}
\begin{itemize}
    \item SaaS-Employee management system (chat, surveys, benchmark, analytics...) and Restful API Development
    \item Socket programming in python to receive continuous data from IoT device and display chunks of data meaningfully
    \item Streaming service with Web-Portal, DRF and Kodi across platforms
\end{itemize}

{\textbf{Freelance Full-stack Developer}}   \hfill {2015 - Apr 2019}
\begin{itemize}
    \item Cloud solution for Data gathering and webscraping with selenium and GCP.
    \item Designing template and banners for High resolution prints
    \item Inventory management system for Book store with Django.
\end{itemize}
% \fi

\end{rSection}


%----------------------------------------------------------------------------------------
%     PROJECT SECTION
%----------------------------------------------------------------------------------------
\begin{rSection}{Projects}

% === 2026 Projects ===
{\bf DevOps Management for Cluster Validation} \hfill {2025 - Present}\\
\textbf{Description}: Delivering consistent CI for cluster validation with S3 log upload and security policy configuration across cloud networks and regulations.

\begin{itemize}
    \item \textbf{Cluster Validation CI:} Developed and delivered consistent Continuous Integration workflows for automated cluster validation.
    \item \textbf{Cloud Infrastructure:} Configured automated log upload processes to AWS S3, ensuring reliable log retention and monitoring.
    \item \textbf{Security \& Compliance:} Implemented appropriate security policy configurations tailored to various cloud network setups and industry regulations.
\end{itemize}

{\bf GPU Performance Analysis for Data Center Workloads (R\&D)} \hfill {2024 - Present}\\
\textbf{Description}: High-Performance Computing research for in-band GPU profiling on distributed systems

\begin{itemize}
    \item \textbf{HPC Profiling Architecture:} Designed and implemented robust in-band profiling mechanisms using \textbf{C++} and \textbf{Python} for comprehensive GPU telemetry across large-scale data center deployments.
    \item \textbf{Distributed Workload Orchestration:} Developed dynamic workload orchestration systems for heterogeneous computing environments (\textbf{CPUs, GPUs, accelerators}) with real-time resource adaptation.
    \item \textbf{Deep Learning Workload Analysis:} Profiled \textbf{TensorFlow} and \textbf{PyTorch} training workloads on multi-GPU clusters, analyzing CUDA/HIP kernel performance and memory bandwidth utilization.
    \item \textbf{Operating System Integration:} Implemented low-level system interfaces for hardware metric collection including GPU temperature, power consumption, and memory management through SMI.
    \item \textbf{Performance Optimization:} Developed kernel profiling tools using \textbf{modern C++} for call stack timing analysis, identifying bottlenecks in high-performance applications.
\end{itemize}

{\bf Intelligent Defect Triage System using RAG \& LangGraph (R\&D)} \hfill {2024 - Present}\\
\textbf{Description}: AI-powered defect analysis and real-time triage prediction for enterprise HPC validation workflows

\begin{itemize}
    \item \textbf{RAG Pipeline Architecture:} Designed and implemented a Retrieval-Augmented Generation (RAG) pipeline using \textbf{LangGraph} and \textbf{Python} for automated JIRA dashboard defect analysis.
    \item \textbf{Defect Trend Analysis:} Developed ML models to identify defect root cause patterns, rejection reason trends, and domain-based cluster classification across HPC validation datasets.
    \item \textbf{Real-Time Triage Prediction:} Built intelligent triage prediction system leveraging contextual availability (setup/infra issues, standard system call failures, known solutions) to accelerate defect resolution.
    \item \textbf{LLM Integration:} Integrated large language models with \textbf{PyTorch} embeddings for semantic similarity search and context-aware defect categorization.
    \item \textbf{Enterprise Scalability:} Deployed solution on large-scale on-prem infrastructure supporting distributed teams across multiple validation clusters.
\end{itemize}

% === 2023 Projects ===
{\bf Prompt Engineering and Generative AI for firmware stack} \hfill {2023 - 2024}\\
\textbf{Description}: Solution design for accelerated silicon tape-in with accelerating firmware deliveries

\begin{itemize}
    \item Work task summarization with Azure LLM Open AI services (GPT3.5 and GPT4)
    \item Code Review against standard static code analysis and common vulnerabilities finding
    \item Generation of Code Enhancement suggestion 
\end{itemize}

\textbf{Platform Provisioning with Pre-Production Reference Board} \hfill {2023 - 2024}\\
\textbf{Description}: Platform Provisioning with automating most common features for pre-production server processor such as:

\begin{itemize}
    \item Remote BMC and BIOS Firmware Flashing
    \item OS Deployment from artifactory through Jenkins Pipeline
    \item Boot to EFI Shell and OS
    \item Configuration of Initial system setup with Jenkins pipeline steps with Baseboard management controller RedFish API
\end{itemize}

% === 2022 Projects ===
\textbf{Universal Scalable Firmware} \hfill {2022 - 2024} \\
\textbf{Description}: Architecture design of Platform orchestration layer for Firmware Configuration

\begin{itemize}
    \item Motivation of workgroup - \href{https://uefi.org/sites/default/files/resources/Firmware%20Configuration%20%E2%80%93%20Past%2C%20Present%2C%20and%20Future_Zimmer.pdf}{Firmware Configuration – Past, Present, and Future}
    \item \href{https://universalscalablefirmware.github.io/documentation/7_yaml_boot_configuration.html}{Firmware Configuration}
\end{itemize}

% === 2020 Projects ===
\textbf{Firmware Automation Framework}  \hfill {2020 - Present} \\
\textbf{Description}: Lead SDL for an automation foundation framework for high scaled volume of silicon devices with edk2 UEFI firmware;

\begin{itemize}
    \item Seamless configuration of system without any manual intervention for BIOS configuration.
    \item SDL activities like legal compliance, security, licensing, static code analysis
    \item Continuous deployment and test with jenkins and tox configurations
    \item Release Artifact management to JFrog Artifactory using GitHub Actions
    \item \href{https://github.com/intel/xml-cli}{https://github.com/intel/xml-cli}
\end{itemize}

% === 2019 Projects ===
\textbf{Multiview Video Summarization} \hfill {2019 - 2020}\\
\textbf{Description}: Summarizing video of multiple view angles of egocentric location in to one to generate effective video summary, minimizing overhead to explore multiple cam feed to discover accidents or anomalies.

% \textbf{Consummating Research Projects using Agile Manifesto} \hfill {2018}\\
% \textbf{Description}: Utilization of Agile Manifesto for consummating research projects. Enlightening the appropriateness of Scrum for the research projects which is a widely used agile manifesto at present in the software industry.

% \textbf{Inventory Management and GST Billing for Mobile Shop} \hfill {2018}\\
% \textbf{Description}: To avoid the headache of manual billing automated GST billing system designed to manage stocks and update according to purchase and sales and generate GST invoice for sales.

% \textbf{Streaming Box} \hfill {2017 - 2018}\\
% \textbf{Description}: An alternate streaming service to DTH and providing easy customized IPTV streaming service based on Kodi (formerly xbmc) Framework

% \begin{itemize}
%     \item Backend Streaming Service, authentication and server security with Django Framework and REST API
%     \item Web Portal to enable vendor customize Streaming Front-end UI for their users
%     \item Cross compilation and Native compilation to enable various platform support as Android, Windows and MacOS
% \end{itemize}

\end{rSection}

%----------------------------------------------------------------------------------------


%----------------------------------------------------------------------------------------
%	EDUCATION SECTION
%----------------------------------------------------------------------------------------

\begin{rSection}{Education}

{\bf M.Tech Computer Science} \hfill {Jul 2018 - Jun 2020 }
\\ 
Nirma University, GPA: 8.49

{\bf B.E. in Computer Science and Engineering} \hfill {Aug 2013 - Apr 2017 }
Gujarat Technological University, CGPA: 8.08

{\textbf{12th Science, GSHSEB, Class XII}}  \hfill {Jun 2011 - Mar 2013 }\\
Science Stream, 74.92\% 
 
% {\textbf{GSHSEB, Class X}}  \hfill {Jun 2010 - Mar 2011 }\\
% 81.4\%

%Minor in Linguistics \smallskip \\
%Member of Eta Kappa Nu \\
%Member of Upsilon Pi Epsilon \\

\end{rSection}

%----------------------------------------------------------------------------------------
%     Certifications
%----------------------------------------------------------------------------------------
\begin{rSection}{Certifications}

\begin{tabular}{ @{} >{\bfseries}l @{\hspace{6ex}} l }
2025 & \href{https://learning.amd.com/certs/11279/95A05A0CC78443498F47FD0FDB6188FC65062.pdf}{AMD ROCm Star - Application Developer} by AMD Adaptive Computing Training Team
\\
2025 & \href{https://learn.deeplearning.ai/accomplishments/f4188f58-c075-45f4-b2c8-6b6318c51ceb?usp=sharing}{On-Device AI} by Qualcomm \& DeepLearning.ai
\\
2025 & \href{https://www.linkedin.com/learning/certificates/bfa1f8149f171120da54a209a31361de4ac56fe1e3f167e2cbc8cd9fd7ca38d1?u=2140730}{Essential Terraform in AWS}
\\
2021 & \href{https://www.credly.com/badges/a3861155-a2e2-43d5-a8dd-cc21651c5841?source=linked_in_profile}{Product Assurance and Security Yellow Belt - Software}
\\
2020 & Natural Language Processing with Classification
and Vector Spaces
\\ & \href{https://coursera.org/verify/S9DGGMCDMBGQ}{Coursera S9DGGMCDMBGQ}
\\
2020 & Improving Deep Neural Networks: Hyperparameter
Tuning, Regularization and Optimization
\\ & \href{https://coursera.org/verify/SMWQJKJYJW6G}{Coursera SMWQJKJYJW6G}
\\
2020 & Structuring Machine Learning Projects - \href{https://coursera.org/verify/SEJML7H3TN5Z}{Coursera SEJML7H3TN5Z}
\\
2020 & Neural Networks and Deep Learning - \href{https://coursera.org/verify/S2TETZHTLUSL}{Coursera S2TETZHTLUSL}
\\
2020 & Machine Learning - \href{https://coursera.org/verify/XVNXNPPZ3JMA}{Coursera XVNXNPPZ3JMA}

\end{tabular}

\end{rSection}

%----------------------------------------------------------------------------------------
%     Achievements
%----------------------------------------------------------------------------------------
\begin{rSection}{Achievements}

\begin{tabular}{ @{} >{\bfseries}l @{\hspace{6ex}} l }
2023 & Firmware Configuration Presentation at UEFI Developers Conference \& Plugfest \\
2023 & Tech Talk on Harnessing Multiprogramming across CPU, GPU and 
\\ & QPU (Quantum Processing Unit) \\
2021 & Expert Lecture on CUDA Programming \\
2019-Present & 4000+ reputation on Stackoverflow \\
2017-Present & Active volunteer on Peer Learning, Cultural and team building activity \\
2018 & Rank \#1 for Python challenges on hackerrank \\
2018 & Rank \#112 from 2 lacs+ participants in Techgig's National Coding Contest 2018 \\
% 2014 & Establishment of complete network for Department

\end{tabular}

\end{rSection}
%----------------------------------------------------------------------------------------


%----------------------------------------------------------------------------------------
%     Publications
%----------------------------------------------------------------------------------------
\begin{rSection}{Publications}

\begin{itemize}
    \item \href{https://youtube.com/playlist?list=PLWyBQeJgIuzDFELUK9ineNkOr9LuYatkR&si=mUVeIkLFK5gzyZa7}{HIP Parallel Programming with ROCm} Lecture Series, YouTube  \hfill {2025}
    \item \href{https://youtube.com/playlist?list=PLWyBQeJgIuzC6oN8AihmSMkgheX9TLDaW&si=HAOmHXKP4KdsAzPX}{Beyond the Power | Data Center \& BMC} Lecture Series, YouTube  \hfill {2025}
    \item \href{https://www.youtube.com/playlist?list=PLWyBQeJgIuzD2o9ZVw5oI-P2fX4uyNKZA}{System Design for Parallel Programming} Lecture Series, YouTube  \hfill {2025}
    % \item A Review on Honeypot : Track the Attack. NCCICT, 2016 \hfill {2016}
    \item \href{https://uefi.org/sites/default/files/resources/Firmware%20Configuration%20%E2%80%93%20Past%2C%20Present%2C%20and%20Future_Zimmer.pdf}{Firmware Configuration – Past, Present, and Future}. Presentation, UEFI DevCon 2023 \hfill {2023}
\end{itemize}


\end{rSection}

\end{document}

